% Bounded research proof. The global metric radius is never differentiated
% in the arbitrary-current test.
\begin{proof}
\emph{The formal geometry and the ordinary-topology distinction.}
The hypersurface in the total space of
\(\mathcal L\oplus\mathcal L^{-1}\) defined by
\[
 (a,b)\longmapsto b(a)=\pi
\]
has local equations \(xy=\pi\) over the smooth curve
\(\mathcal C\). Its completion along the special fibre is a
strictly semistable formal surface. Its two components are
the total spaces of \(\mathcal L_0\) and \(\mathcal L_0^{-1}\),
meeting along their zero sections \(C_0\). The two normal
bundles of this double curve are respectively
\(\mathcal L_0\) and \(\mathcal L_0^{-1}\).
On the generic fibre the vector \(a\) is nonzero and determines
\(b\), so the open locus with \(0<u<\ell\) is precisely \(X\).
In particular \(X\) is smooth, hence quasi-smooth, of
dimension two and is without analytic boundary.

The formal metric gives a continuous norm on the analytic
total space, by \cite[Sections 17.1.2, 17.1.4--17.1.5 and
17.7.2, pp.~471--473 and 493]{CLD25}. Its local formal frames
are orthonormal. Formal Zariski trivializations give analytic
domains, not necessarily ordinary open neighbourhoods. We
therefore use them neither as an ordinary open cover nor to
differentiate \(u\) globally. This is the distinction between
ordinary forms and G-forms in
\cite[Proposition 15.1.1, pp.~421--422]{CLD25}.

We will use the following elementary compactness fact. For
\(0<a\leq b<\ell\), the strip
\begin{equation}
\label{eq:p213-torsor-compact-strip}
 \{\xi\in L^\times:a\leq u(\xi)\leq b\}
 \quad\text{is compact.}
\end{equation}
Indeed the total line bundle has its projective-bundle
compactification by adjoining the infinity section. This
compactification is proper over the proper curve \(C\), so
its analytification is compact. The norm is continuous on
the finite chart including the zero section and vanishes
there. On the chart at infinity, the inverse vector has the
dual norm, continuous and vanishing at the infinity section.
The displayed lower and upper norm bounds therefore define
a closed subset of that compactification disjoint from both
sections. No assertion that \(X\) itself is compact is made.

\emph{Each residual point still gives an actual product collar.}
Perfection of \(k\), henselian lifting, and smoothness give
the finite unramified \(R_p/R\) and the completed local ring
\(R_p[[w_p]]\). The finite extension has degree \(d_p\).
The local ring is regular of dimension two with maximal
ideal generated by \(\pi,w_p\); successive approximation
gives the power-series presentation. The tube is the open
unit disc over \(K_p\). One may verify this on the subdomains
\(\{|w_p|^N\leq|\pi|\}\), which exhaust the tube.
A basis \(e_p\) of the completed line bundle gives an analytic
basis on this entire disc, of formal norm one. These assertions
are the same completed-ring and convergence calculations as
in the separately audited product packet, now applied also
to a rank-one free module.

Thus the inverse image of an outer disc collar is the actual
ordinary open \(K\)-analytic domain
\[
 V_p=\{|\pi|<|x_p|<1\}_{K_p}
          \times_{K_p}\{e^{-\delta_p}<|w_p|<1\},
 \qquad q_p=(v(x_p),v(w_p)).
\]
The inverse image under \(q_p\) of a compact subset of \(Q_p\)
is closed in a compact closed polyannulus. Pullbacks of
compactly supported target forms, including all primitives
used below, consequently have compact analytic support.
Their extension by zero to \(X\) is legitimate. This is the
actual support requirement in
\cite[Section 7.1.11, p.~159]{CLD25}.

For completeness, the collar coefficient requires no
assumption about a representation of \(T\). Pairing \(T\)
with \(q_p^*(h\Omega_p)\) gives a positive distribution
\(\nu_p\) on \(Q_p\). A nonnegative smooth cutoff equal to
one on any prescribed compact subset gives the sup-norm
bound there, so \(\nu_p\) is a locally finite positive
Radon measure. Deleting \(\mathrm d'u\), respectively
\(\mathrm d's\), from \(\Omega_p\) with the appropriate
sign and using \(\mathrm d'T=0\) shows that
\[
 \nu_p(\partial_u h)=\nu_p(\partial_s h)=0.
\]
These primitives have compact analytic support by the
preceding paragraph. On any relatively compact target ball,
Euclidean convolution has both derivatives zero and is
constant. Testing with a function of integral one shows
that the constants converge to a finite constant as the
mollifier shrinks. Overlapping balls and the connectedness
of \(Q_p\) give \(\nu_p=b_T(p)\,\mathrm du\,\mathrm ds\),
with \(0\leq b_T(p)<\infty\).
This also proves independence under shrinking the collar.

If \(e'_p=g e_p\) is another completed formal basis,
\(g\in R_p[[w_p]]^\times\). Its constant coefficient is
a unit and all higher coefficients are bounded, so
\(|g|=1\) on the entire tube. Therefore \(v(x'_p)=v(x_p)\)
as actual functions. For another formal parameter
\(w'_p=a_0+a_1w_p+\cdots\), one has
\(a_0\in\pi R_p\), \(a_1\in R_p^\times\) and
\(a_j\in R_p\). On a sufficiently outer collar the linear
term strictly dominates all others, giving
\(|w'_p|=|w_p|\). Both independence assertions compare
actual forms on an ordinary open common collar, not just
values on a skeleton. The coefficient-algebra identification
is fixed by its residue map through henselian uniqueness.

\emph{A rational section supplies an ordinary-open coordinate.}
The model is regular with one integral special fibre. At
its generic point \(\eta_0\), the line bundle is free over
the local DVR. Choose a basis there and view it as a
nonzero rational section \(\sigma\) of \(\mathcal L\).
Its vertical divisor coefficient is zero, and its reduction
is a nonzero rational section \(\bar\sigma\) of
\(\mathcal L_0\). Write
\[
 n_p=\operatorname{ord}_p(\bar\sigma).
\]
This does not trivialize \(\mathcal L_0\), and the divisor
\(\sum n_pp\) need not have degree zero.

Let \(f\in k(C_0)^\times\). Lift it to a unit \(F\) of the
same generic-special-fibre DVR. Then \(F\) is a rational
function on \(C\) with zero vertical order; put
\(m_p=\operatorname{ord}_p(f)\). Let \(B\) be the union of
all specializations of zeros and poles on \(C\) of both
\(\sigma\) and \(F\). It is finite: each horizontal divisor
has finite intersection with the special fibre of this
proper regular surface. Points with \(m_p=0\), or with
\(n_p=0\), can still belong to \(B\) and are retained.

In the completed frame write \(\sigma=a_pe_p\).
The meromorphic coefficient \(a_p\) and the function \(F\)
are quotients of elements of \(R_p[[w_p]]\) whose vertical
orders can be cancelled. The first nonzero reduced
coefficient of each numerator and denominator is a unit.
On a sufficiently outer collar its monomial strictly
dominates all lower terms, whose coefficients lie in
\(\pi R_p\), and all higher bounded terms. Hence on the
whole ordinary collar one has the exact identities
\begin{equation}
\label{eq:p213-torsor-two-slopes}
 h_\sigma:=-\log\|\sigma\|_{\mathcal L}=n_ps_p,
 \qquad t:=v(F)=m_ps_p .
\end{equation}
Outside the tubes of \(B\), both \(\sigma\) and \(F\) are
formal units in their respective local frames. Their
norms are one there, so \(h_\sigma=t=0\).
This last assertion is a pointwise assertion on that
set; it does not assert ordinary-open constancy at the
good-reduction point.

Remove from \(C^{\mathrm{an}}\) the finite support of
\(\operatorname{div}_C(\sigma)\) and
\(\operatorname{div}_C(F)\), and denote the resulting
ordinary open set by \(U\). Over \(U\), \(\sigma\) is
an invertible analytic section, not merely a formal
G-local frame. For \(\xi\in L^\times|_U\), write
\(\xi=a\,\sigma(\rho(\xi))\). The coefficient \(a\)
is a globally defined invertible analytic function
on this ordinary open total space. Thus
\[
 r=v(a),\quad t=v(F)\circ\rho
 \quad\text{are ordinary smooth log-unit functions,}
 \qquad u=r+h_\sigma\circ\rho.
\]
In particular \(\mathrm d'\mathrm d''r
=\mathrm d'\mathrm d''t=0\) there. On the collar,
\begin{equation}
\label{eq:p213-torsor-shear}
 r=u-n_ps,\quad t=m_ps,\qquad
 \begin{pmatrix}r\\s\end{pmatrix}
 =\begin{pmatrix}1&-n_p\\0&1\end{pmatrix}
   \begin{pmatrix}u\\s\end{pmatrix}.
\end{equation}
The integral lattice shear has determinant one.
It is not a globally defined annulus coordinate, and
it does not imply that \(u\) is globally smooth.

\emph{Uniform smallness and compact support in the torsor.}
Choose \(\varphi\in C_c^\infty((0,\ell))\) with integral
one, and numbers \(0<A<B'<\ell\) containing its support.
Fix
\[
 0<\gamma<\min(A,\ell-B').
\]
Shrink the collars in \eqref{eq:p213-torsor-two-slopes}
and choose, for every \(p\in B\),
\[
 0<a_p'<b_p'<\delta_p,\qquad |n_p|b_p'<\gamma.
\]
The condition is empty when \(n_p=0\).
Choose a smooth cutoff \(\eta_p(s)\), equal to one for
\(s\leq a_p'\), zero for \(s\geq b_p'\), and between zero
and one. Define \(\eta\) on the base curve to be
\(\eta_p(s_p)\) on the tube of \(p\in B\), zero at its
centre, and one elsewhere.
The function \(1-\eta_p\) is constant near the centre
and supported in a compact closed disc strictly
inside the ordinary open tube. Its extension by zero
outside that tube is smooth. There are only finitely
many such functions, proving that \(\eta\) is globally
ordinary CLD smooth.

The cutoff vanishes on neighbourhoods of every zero
and pole of both \(\sigma\) and \(F\), since the exact
outer-collar identities put these divisors in the
deeper region. The base is proper, so
\(\operatorname{supp}(\eta)\) is compact and contained
in \(U\). The continuous function \(h_\sigma\) is
defined on \(U\), and the construction gives
\begin{equation}
\label{eq:p213-torsor-smallness}
 |h_\sigma|\leq\max_{p\in B}|n_p|b_p'<\gamma
       \quad\text{on }\operatorname{supp}(\eta).
\end{equation}
If \(B\) is empty, the left side is zero and no maximum
is needed.

Assume that some \(m_p\neq0\), since otherwise the
principal relation is trivial. Choose
\[
 0<\varepsilon<
     \min_{\substack{p\in B\\m_p\neq0}}|m_p|a_p'
\]
and \(\chi\in C_c^\infty((-\varepsilon,\varepsilon))\)
equal to one near zero. On \(X\cap L^\times|_U\), use
the smooth \((1,2)\)-form
\[
 \beta=\varphi(r)\chi(t)\,
             \mathrm d'r\wedge\mathrm d''r\wedge\mathrm d''t.
\]
The form \((\eta\circ\rho)\beta\), extended by zero,
is a compactly supported ordinary smooth form on \(X\).
Indeed the removed base divisor has a neighbourhood
where \(\eta=0\), so extension by zero is smooth.
On its support, \(r\in[A,B']\) and
\eqref{eq:p213-torsor-smallness} gives
\[
 A-\gamma\leq u\leq B'+\gamma,
 \qquad 0<A-\gamma<B'+\gamma<\ell.
\]
It follows from \eqref{eq:p213-torsor-compact-strip}
that the support is compact inside \(X\).
This is the radial support estimate that would fail
if one simply substituted an arbitrary rational
coordinate for the global product coordinate.

\emph{The cutoff error and the complete principal-divisor test.}
For notation suppress pullback by \(\rho\).
The support of \(\mathrm d'\eta\) lies in the closed
transition collars \(a_p'\leq s\leq b_p'\).
If \(m_p\neq0\), then \(|t|\geq|m_p|a_p'>\varepsilon\),
so \(\chi(t)=0\) there. If \(m_p=0\), then
\eqref{eq:p213-torsor-two-slopes} gives \(t=0\) on an
entire ordinary collar, so \(\mathrm d''t=0\).
This includes every case with \(m_p=0\) and
\(n_p\neq0\), as well as deeper cancellations with
\(m_p=n_p=0\). Thus, as source forms before pairing,
\[
 \mathrm d'\eta\wedge\beta=0.
\]
Since \(r,t\) are actual log-units where the form lives,
direct differentiation gives
\[
 \mathrm d'\beta
   =\varphi(r)\chi'(t)\,
       \mathrm d'r\wedge\mathrm d''r
                       \wedge\mathrm d't\wedge\mathrm d''t.
\]
No derivative of the continuous global \(u\) or of a
formal metric appears in this computation.

Where \(\eta\neq0\), the support of this last form is
contained in the finitely many individual collars
with \(m_p\neq0\). Outside the \(B\)-tubes, \(t=0\);
continuity and the fact that \(\chi'\) vanishes on an
open interval around zero make the form vanish on an
ordinary neighbourhood of every such point. This
also handles the good-reduction point without any
ordinary-open interpretation of a formal Zariski
cover. In the \(p\)-tube the support has
\(0<s<\varepsilon/|m_p|<a_p'\), hence \(\eta=1\).
It is bounded away from \(s=0\), since \(\chi'\)
vanishes near zero. The radial margin above bounds
\(u\) away from both ends. These are compact supports
inside the disjoint ordinary source collars.

By \eqref{eq:p213-torsor-shear}, on such a collar
\begin{equation}
\label{eq:p213-torsor-top-shear}
 \mathrm d'\beta=
 m_p^2\varphi(u-n_ps)\chi'(m_ps)\,
    \mathrm d'u\wedge\mathrm d''u\wedge
                     \mathrm d's\wedge\mathrm d''s.
\end{equation}
Terms containing an additional \(\mathrm d's\) or
\(\mathrm d''s\) vanish. There is no additional
factor from \(n_p\), because the shear has determinant
one. For every \(s\) relevant to this support,
\(|n_p|s<\gamma\), and consequently
\[
 \int_0^\ell\varphi(u-n_ps)\,\mathrm du
       =\int_{\mathbf R}\varphi(r)\,\mathrm dr=1.
\]
The smooth source form \(\eta\,\mathrm d'\beta\)
is the finite sum of the zero extensions of these
collar forms. Closedness and \eqref{eq:p213-torsor-collar}
therefore imply
\[
 \begin{split}
 0&=\langle T,\mathrm d'(\eta\beta)\rangle\\
  &=\sum_{\substack{p\in B\\m_p\neq0}} b_T(p)
       \int_0^{\delta_p}m_p^2\chi'(m_ps)
           \left(\int_0^\ell\varphi(u-n_ps)\,\mathrm du\right)
                      \mathrm ds\\
  &=-\sum_p m_pb_T(p).
 \end{split}
\]
The final sign follows, for either sign of \(m_p\), from
\[
 \int_0^{\delta_p}m_p^2\chi'(m_ps)\,\mathrm ds
   =m_p\bigl(\chi(m_p\delta_p)-\chi(0)\bigr)=-m_p.
\]
The inner integral equals one wherever the outer
integrand is nonzero; no assertion is required at
other values of \(s\). This proves
\eqref{eq:p213-torsor-principal} for every \(f\).
It also shows precisely why the normal-bundle orders
\(n_p\) cause no curvature correction in this particular
principal-divisor relation.

\emph{Positivity, degree, and the conditional calibration reading.}
The finite numbers \(b_T(p)\) define a finite-valued
homomorphism on the divisor group, nonnegative on
effective divisors. Here finite-valued does not mean
a finite total measure on the set of all closed points.
For a degree-zero divisor \(D\), choose an effective
\(H\) with \(\deg H+\chi(\mathcal O_{C_0})>0\).
Riemann--Roch gives a nonzero section of
\(\mathcal O(H+jD)\) for every integer \(j\).
Positivity and the principal relation give
\(0\leq\lambda_T(H)+j\lambda_T(D)\) for both signs
of \(j\), so \(\lambda_T(D)=0\).
Applying this to \((\deg q)p-(\deg p)q\) gives
\(b_T(p)=c\,d_p\) for a single \(c\geq0\).

As a \(K_p\)-analytic product polyannulus the collar
has canonical coefficient one in primitive coordinates.
When the same space is viewed over \(K\), that coefficient
is multiplied by \([K_p:K]=d_p\), by
\cite[Remark 14.4.3, p.~410]{CLD25}.
The extension is unramified; no lattice ramification
factor is present. There is a single source Gauss plane,
not \(d_p\) ordinary sheets. This is the same explicit
finite-constant-field source as in the product packet,
and the integral shear does not change its calibration.

Under the separately assumed source-calibration
representation, compact tests on this injectively
tropicalized source Gauss plane identify its coefficient
with \(b_T(p)\), hence its ratio to the canonical
coefficient with \(c\). The geometric and current tests
above do not use that representation.
No passage to other source cells or to equality
\(T=c\delta_X\) on \(X\) is justified by this conclusion.
\end{proof}

